\documentclass[11pt,a4paper]{article}

\usepackage[a4paper,margin=2.5cm]{geometry}
\usepackage{fontspec}
\setmainfont{Roboto}[
  Path = fonts/,
  Extension = .ttf,
  UprightFont = *-Regular,
  BoldFont = *-Bold,
  ItalicFont = *-RegularItalic,
  BoldItalicFont = *-BoldItalic]
\setmonofont{RobotoMono}[Path = fonts/, Extension = .ttf, UprightFont = *-Regular, Scale = 0.9]
\usepackage{polyglossia}
\setdefaultlanguage{french}
\setotherlanguage{english}
\usepackage{float}
\usepackage{array}
\usepackage{microtype}
\usepackage{titlesec}
\usepackage[square,authoryear]{natbib}
\usepackage{xcolor}
\usepackage{booktabs}
\usepackage{tabularx}
\usepackage{tikz}
\usetikzlibrary{arrows.meta,positioning}
\usepackage{caption}
\usepackage[hidelinks=false,colorlinks=true,linkcolor=black,citecolor=citeblue,urlcolor=citeblue]{hyperref}

\definecolor{citeblue}{HTML}{1F4E79}

\linespread{1.12}
\setlength{\emergencystretch}{3em}
\setlength{\parindent}{1.2em}
\setlength{\parskip}{0pt}
\titleformat{\section}{\large\bfseries}{}{0pt}{}
\titlespacing*{\section}{0pt}{14pt}{6pt}
\captionsetup{font=small,labelsep=endash}
\gappto\captionsfrench{\renewcommand{\figurename}{Figure}\renewcommand{\tablename}{Tableau}\renewcommand{\refname}{Références}}
\setcitestyle{citesep={,}}
\pagestyle{plain}

\begin{document}

\begin{center}
{\Large\bfseries OpenCode, l'alternative open source à Claude Code\\[2pt]
\large Comprendre et utiliser les agents IA pour développer\par}
\vspace{10pt}
Projet FYC de [Prénom Nom], [Prénom Nom], [Prénom Nom] et [Prénom Nom]\\
encadré par [Nom du mentor]\\
ESGI, 4\textsuperscript{e} année, 2026-2027\\
(Fiche scientifique préparatoire)
\end{center}

\section*{Contexte et état de l'art}

Depuis quelques années, les outils d'aide au développement fondés sur l'intelligence artificielle se sont largement diffusés. Les premiers outils d'autocomplétion se limitaient à terminer la ligne en cours. L'arrivée des grands modèles de langage (LLM), construits sur l'architecture Transformer \citep{vaswani2017}, a permis de générer des fonctions entières à partir du contexte d'un fichier, comme le fait GitHub Copilot à partir du modèle Codex \citep{chen2021}. Le réglage de ces modèles par retour humain \citep{ouyang2022} a ensuite rendu possible un usage conversationnel, popularisé par ChatGPT fin 2022.

Ces usages restent cependant passifs : le modèle propose, et c'est le développeur qui exécute, teste et recopie le code. Les travaux sur les agents ont montré qu'un modèle peut alterner raisonnement et actions, en appelant des outils puis en observant leur résultat \citep{yao2023,schick2023}. C'est sur ce principe que reposent les agents de développement apparus en 2025, comme Claude Code \citep{anthropic2025} : un modèle, un ensemble d'outils (lecture et écriture de fichiers, terminal, recherche dans le code) et une boucle qui s'exécute jusqu'à la fin de la tâche. Leur capacité à résoudre de vrais tickets issus de projets open source est aujourd'hui mesurée par des bancs d'essai comme SWE-bench \citep{jimenez2024}.

L'efficacité d'un agent dépend fortement du contexte qu'il reçoit. Les modèles exploitent moins bien une information placée au milieu d'un long contexte \citep{liu2024}, ce qui pousse à sélectionner les informations utiles plutôt qu'à tout fournir. Pour apporter des données externes au modèle, deux approches coexistent : la génération augmentée par récupération, ou RAG \citep{lewis2020}, et le Model Context Protocol (MCP), un standard ouvert qui permet à l'agent d'accéder lui-même à des outils et à des sources de données \citep{anthropic2024}. Le serveur MCP Context7 \citep{context7} en est un exemple : il fournit à l'agent la documentation à jour des bibliothèques utilisées dans le projet.

Les agents de développement les plus utilisés sont pour la plupart propriétaires et liés à un seul fournisseur de modèles. OpenCode \citep{opencode} propose une alternative : un agent en terminal publié sous licence MIT, compatible avec de nombreux fournisseurs de modèles ainsi qu'avec des modèles exécutés en local.

\section*{Sujet et problématique}

Notre projet porte sur OpenCode et, plus largement, sur l'usage des agents IA dans le développement logiciel. Il prend la forme d'un cours vidéo destiné à des étudiants en développement web de niveau Mastère.

Le choix d'un agent ne se limite pas à la qualité du code qu'il produit. Il engage aussi le coût d'utilisation, facturé au nombre de tokens traités, la confidentialité du code envoyé au fournisseur et la dépendance envers un éditeur. Un outil fermé comme Claude Code bénéficie d'une intégration poussée avec les modèles de son éditeur. Un outil ouvert laisse le choix du modèle, mais ses résultats dépendent alors de ce choix. Enfin, mal utilisé, un agent peut produire du code faux ou exécuter des actions non souhaitées.

Notre problématique est donc la suivante : OpenCode est-il une alternative crédible à Claude Code, et comment l'utiliser efficacement et en sécurité sur un projet réel ?

\section*{Méthode}

Nous abordons cette question en trois temps. D'abord, nous construisons les notions nécessaires pour comprendre ce qu'est un agent. Le cours part du fonctionnement d'un LLM (prédiction du token suivant, entraînement, limites), puis présente la boucle agentique (figure~\ref{fig:boucle}), le rôle du contexte et le protocole MCP. Nous présentons ensuite Claude Code et OpenCode, ainsi que l'architecture client-serveur de ce dernier (figure~\ref{fig:archi}).

Nous passons ensuite à la pratique sur un projet de taille modeste : une API Node.js/Express accompagnée de tests. L'apprenant installe et configure OpenCode, puis réalise un cas pratique en trois parties : comprendre un projet qu'il ne connaît pas, y ajouter une fonctionnalité (la pagination d'un endpoint), puis corriger un bug signalé par un test. Chaque partie est accompagnée d'une correction en vidéo.

Enfin, nous comparons les deux outils à partir de leur documentation officielle, sur des critères qui comptent pour un développeur ou une entreprise : licence, modèles utilisables, coût, devenir des données et maturité (tableau~\ref{tab:comparatif}). Le cours se termine par des bonnes pratiques d'usage des agents, notamment sur la relecture du code produit et la sécurité.

\section*{Résultats}

Le cours se compose de douze vidéos courtes, pour une durée totale d'environ 1\,h\,35, et d'un cas pratique en trois parties. À la fin du cours, l'apprenant doit être capable d'expliquer le fonctionnement d'un agent IA, d'installer et de configurer OpenCode avec le fournisseur de son choix, puis de l'utiliser sur un projet existant en gardant la main sur les modifications proposées.

La comparaison montre qu'OpenCode est moins abouti que Claude Code, dont l'intégration avec les modèles d'Anthropic est très poussée. En revanche, il laisse au développeur le choix du modèle, la maîtrise des coûts et la possibilité de garder ses données en local. Il constitue donc une alternative crédible lorsque ces critères comptent, à condition de l'utiliser avec méthode.

\section*{Aller plus loin}

OpenCode évolue rapidement, et certaines commandes présentées dans le cours peuvent changer d'une version à l'autre. La qualité des résultats dépend aussi beaucoup du modèle branché sur OpenCode.

L'autonomie d'un agent pose enfin des questions de sécurité. Un agent qui lit des fichiers, des pages web ou des tickets peut y trouver des instructions malveillantes, une attaque appelée injection de prompt indirecte \citep{greshake2023}. Un prolongement de ce travail serait d'étudier les mécanismes de permission et d'isolation proposés par ces outils, et la façon de les configurer dans un contexte d'entreprise.

\begin{thebibliography}{99}
\small
\begin{english}

\bibitem[Anthropic, 2024]{anthropic2024}
Anthropic (2024). Introducing the Model Context Protocol. \url{https://www.anthropic.com/news/model-context-protocol}

\bibitem[Anthropic, 2025]{anthropic2025}
Anthropic (2025). Claude Code overview. Documentation officielle. \url{https://docs.claude.com/en/docs/claude-code/overview}

\bibitem[Chen et~al., 2021]{chen2021}
Chen, M., Tworek, J., Jun, H., Yuan, Q., et al. (2021). Evaluating Large Language Models Trained on Code. arXiv:2107.03374.

\bibitem[Greshake et~al., 2023]{greshake2023}
Greshake, K., Abdelnabi, S., Mishra, S., Endres, C., Holz, T., and Fritz, M. (2023). Not what you've signed up for: Compromising Real-World LLM-Integrated Applications with Indirect Prompt Injection. In \emph{AISec '23}, ACM.

\bibitem[Jimenez et~al., 2024]{jimenez2024}
Jimenez, C. E., Yang, J., Wettig, A., Yao, S., Pei, K., Press, O., and Narasimhan, K. (2024). SWE-bench: Can Language Models Resolve Real-World GitHub Issues? In \emph{ICLR}.

\bibitem[Lewis et~al., 2020]{lewis2020}
Lewis, P., Perez, E., Piktus, A., et al. (2020). Retrieval-Augmented Generation for Knowledge-Intensive NLP Tasks. In \emph{NeurIPS}.

\bibitem[Liu et~al., 2024]{liu2024}
Liu, N. F., Lin, K., Hewitt, J., Paranjape, A., Bevilacqua, M., Petroni, F., and Liang, P. (2024). Lost in the Middle: How Language Models Use Long Contexts. \emph{Transactions of the ACL}, 12:157--173.

\bibitem[Ouyang et~al., 2022]{ouyang2022}
Ouyang, L., Wu, J., Jiang, X., et al. (2022). Training language models to follow instructions with human feedback. In \emph{NeurIPS}.

\bibitem[Schick et~al., 2023]{schick2023}
Schick, T., Dwivedi-Yu, J., Dessì, R., et al. (2023). Toolformer: Language Models Can Teach Themselves to Use Tools. In \emph{NeurIPS}.

\bibitem[SST, 2025]{opencode}
SST (2025). OpenCode, documentation officielle. \url{https://opencode.ai/docs}

\bibitem[Upstash, 2025]{context7}
Upstash (2025). Context7. \url{https://context7.com}

\bibitem[Vaswani et~al., 2017]{vaswani2017}
Vaswani, A., Shazeer, N., Parmar, N., et al. (2017). Attention Is All You Need. In \emph{NeurIPS}.

\bibitem[Yao et~al., 2023]{yao2023}
Yao, S., Zhao, J., Yu, D., Du, N., Shafran, I., Narasimhan, K., and Cao, Y. (2023). ReAct: Synergizing Reasoning and Acting in Language Models. In \emph{ICLR}.

\end{english}
\end{thebibliography}

\clearpage
\section*{Annexes (figures et tableaux)}

\begin{figure}[H]
\centering
\begin{tikzpicture}[
  box/.style={draw, rounded corners=2pt, minimum width=3.4cm, minimum height=1.1cm, align=center, font=\small},
  arr/.style={-{Stealth[length=2.2mm]}, semithick}]
\node[box] (p) at (0,2.2) {Perception\\{\footnotesize le contexte est envoyé}};
\node[box] (d) at (4.6,0) {Décision\\{\footnotesize le LLM choisit un outil}};
\node[box] (a) at (0,-2.2) {Action\\{\footnotesize l'outil est exécuté}};
\node[box] (o) at (-4.6,0) {Observation\\{\footnotesize le résultat est ajouté}};
\node[font=\small, align=center, text=black!70] at (0,0) {boucle répétée\\jusqu'à la fin de la tâche};
\draw[arr] (p.east) to[bend left=20] (d.north);
\draw[arr] (d.south) to[bend left=20] (a.east);
\draw[arr] (a.west) to[bend left=20] (o.south);
\draw[arr] (o.north) to[bend left=20] (p.west);
\end{tikzpicture}
\caption{Fonctionnement d'un agent IA : la boucle perception, décision, action, observation}
\label{fig:boucle}
\end{figure}

\vspace{6pt}

\begin{figure}[H]
\centering
\begin{tikzpicture}[
  box/.style={draw, rounded corners=2pt, minimum height=1cm, align=center, font=\small},
  arr/.style={-{Stealth[length=2.2mm]}, semithick}]
\node[box, minimum width=3cm] (tui) at (0,0) {Interface terminal\\{\footnotesize (client)}};
\node[box, minimum width=3.6cm] (srv) at (4.6,0) {Serveur OpenCode\\{\footnotesize boucle, outils, sessions}};
\node[box, minimum width=4cm] (cloud) at (10,1.5) {Fournisseurs par API\\{\footnotesize Anthropic, OpenAI, Google\ldots}};
\node[box, minimum width=4cm] (local) at (10,0) {Modèle local\\{\footnotesize Ollama}};
\node[box, minimum width=4cm] (mcp) at (10,-1.5) {Serveurs MCP\\{\footnotesize Context7\ldots}};
\node[box, minimum width=3.6cm] (proj) at (4.6,-2.1) {Fichiers du projet\\{\footnotesize dont AGENTS.md}};
\draw[arr, <->] (tui) -- (srv);
\draw[arr, <->] (srv.east) -- (cloud.west);
\draw[arr, <->] (srv.east) -- (local.west);
\draw[arr, <->] (srv.east) -- (mcp.west);
\draw[arr, <->] (srv) -- (proj);
\end{tikzpicture}
\caption{Architecture d'OpenCode : un serveur local relié au modèle choisi, aux outils et au projet}
\label{fig:archi}
\end{figure}

\vspace{6pt}

\begin{table}[H]
\centering
\small
\captionsetup{position=top}
\caption{Critères de comparaison entre Claude Code et OpenCode}
\label{tab:comparatif}
\begin{tabularx}{0.9\textwidth}{l >{\raggedright\arraybackslash}X >{\raggedright\arraybackslash}X}
\toprule
Critère & Claude Code & OpenCode \\
\midrule
Licence & Propriétaire & Open source (MIT) \\
Modèles & Modèles Anthropic uniquement & Nombreux fournisseurs, modèles locaux \\
Coût & Abonnement ou facturation à l'usage & Outil gratuit, modèle payé au fournisseur \\
Données & Envoyées à Anthropic & Fournisseur au choix, ou traitement local \\
Consignes projet & Fichier CLAUDE.md & Fichier AGENTS.md (format ouvert) \\
Maturité & Très aboutie & Récente, évolue rapidement \\
\bottomrule
\end{tabularx}
\end{table}

\end{document}
